\section{Production Gap}

The theory in this paper is deliberately cleaner than a trading system. That is the right
choice for the current research question, but it creates a production gap.

\paragraph{Local approximation.}
The Taylor expansion is local. It is best for short horizons and fails around jumps,
discrete dividends, exercise-boundary changes, and large volatility-surface repricings.
The residual covariance term absorbs part of this misspecification but does not eliminate
it.

\paragraph{Rolling bucket returns.}
The empirical option assets are rolling buckets, not permanent listed contracts. This is
necessary because listed contracts expire. The bucket panel checks whether option-style
exposures can be allocated with Markowitz machinery; it is not a fill-level backtest. The
current runner marks short-dated selected options to listed-expiry payoff using raw daily
underlying closes and split conversion. That is a stricter holding-period test than a
monthly bucket-to-bucket return, but it is still not a fill-level option backtest with
tradable liquidation quotes, margin, exercise, and assignment handling.

\paragraph{Expected returns.}
The training mean of option-bucket returns is a noisy estimate. The paper's strong
test-window variation should therefore be read as evidence about the measurement stack,
not as evidence of a stable deployable premium.

\paragraph{No implementation costs by design.}
Transaction costs, slippage, margin financing, borrow frictions, and assignment frictions
are excluded because the object is a theory paper. The option premium itself is not
excluded: empirical P\&L is mark-to-market P\&L on paid or sold option premium. The
excluded frictions belong in a separate implementation-cost layer, not in the definition
of the option-only Markowitz analogue.

\paragraph{American single-name options.}
The data layer uses American-model/de-Americanized Greek proxies. This is necessary
because listed U.S. single-name equity options are American style, while the local Taylor
map uses European-equivalent sensitivities as a standardized risk representation. A
production implementation would need a full surface-construction and exercise-policy
module before capital could be deployed.

\paragraph{Capacity and constraints.}
The constraints in the paper are mathematical risk budgets: gross NAV, short-premium,
contract concentration, underlying concentration, delta, gamma, and vega. They are not
broker margin rules, market-impact limits, or locate constraints. A production book would
have to replace or augment them with executable order size, spread, open interest,
assignment, capital, and stress-margin constraints.

These gaps do not weaken the theory. They define what the theory has and has not solved.
The paper supplies the option-only mean--variance object. A production system would supply
the trading, financing, and capacity layer.
